\section{Experimental Setup}
\label{sec:setup}
\paragraph{Write-stage identification.}
We use released WebArena Shopping trajectories from the ReasoningBank/fragility substrate. The original intervention requests six outcome-balanced source trajectories; four return complete memory outputs under both reward branches and two failure-branch calls terminate before assistant text. We do not impute or re-POST these missing arms. A disjoint eight-trajectory prompt control compares released reward-mode divergence with a larger semantic-preserving rewording within each mode. For breadth, 16 additional sources are deterministically frozen before new writer outputs, balanced 8/8 by original outcome and spanning 11 intent-template families. A uniform 4,096-token execution repair regenerates all 32 breadth writer units fresh after the lower-cap parent reveals length censoring. Combined with the original complete sample, this yields 20 paired sources.

\paragraph{Four transport gates.}
We evaluate state construction, exposure, branch-specific uptake, and terminal outcome separately. Forced injection fully crosses the four original complete memory pairs with four frozen future evidence packets; this bypasses native retrieval and tests latent leverage. Exposure is measured by exact replay of the released all-MiniLM-L6-v2 task-description retriever with top-$1$ and threshold 0.3. After source expansion, all held-out Shopping retrieval hits are screened only for released evidence and deterministic evaluation, producing 36 native tasks before terminal outcomes. Each is bound to the source selected by the retriever itself.

\paragraph{Native outcome and representation controls.}
On the 36-task support, success- and failure-conditioned memories use exact stored bytes and the released ReasoningBank human-memory wrapper; future evidence, Doubao Seed 2.0 Mini policy, temperature, evaluator, and four-rollout depth are matched. A no-memory arm uses the same support. A raw writer-input arm substitutes the deterministic compact action/result trace supplied to the memory writer before reward conditioning. Terminal practical-effect tests retain the frozen 0.15 floor together with permutation $p<0.05$.

The strongest same-information representation control is an outcome-blind structured writer. For all 20 sources, DeepSeek-V4-Flash receives the same task prompt and exact action-summary bytes used by the reward-conditioned evidence, at temperature zero and a 4,096-token cap. The control prompt uses the same at-most-three Title/Description/Content schema but contains no success/failure, reward, score, or outcome semantics. All 20 writer calls are frozen before execution; no failed unit may be retried, replaced, or imputed. The 11 neutral memories needed by native retrieval are then serialized with the same wrapper. On the same 36 tasks, 144 fresh terminal calls estimate Eq.~\ref{eq:neutral-control}; both the 0.15 magnitude floor and $p<0.05$ are required.

\paragraph{Pre-terminal uptake.}
To localize attenuation before terminal scoring, the first-action probe freezes released trajectory step 1 for all 36 native tasks, removes any pre-existing task-history memory from the state message, and compares success-memory, failure-memory, and no-memory first structured actions with four rollouts each. The preregistered statistic is mean per-state empirical S/F total-variation distance, using the inherited 0.20 process floor and a 100,000-draw within-state permutation test. No-memory geometry and later next-goal text are descriptive only.

\paragraph{Reliability and scope.}
All scientific runs are inference-only: no training and no local GPU fine-tuning. Provider responses are archived before analysis where supported; each outcome-blind writer/terminal provider POST additionally writes a resumable stage JSON and refreshes the aggregate result. Missing scientific units are never replaced by outcome-favorable tasks. A DeepSeek-V4-Flash downstream-policy transfer stops after no-text length censoring on the first frozen unit at both allowed output caps and therefore contributes zero scientific units rather than a null result. Appendix~\ref{sec:execution-accounting} reports execution accounting and separates scientific non-passes from support/runtime failures.
